%_____________________________________________________________________________________________ 
% Chapter 6: Discussion
%_____________________________________________________________________________________________ 
\chapter{Discussion}
\label{sec:discussion}

\section{Design Trade-offs}

\begin{itemize}
  \item \textbf{In-memory block tree.}  We keep the entire block tree in RAM, which is fine
    for thousands of blocks but will not scale to millions.  A natural extension is to page
    older blocks to disk, using the WAL infrastructure that already exists.

  \item \textbf{Consensus-driven membership.}  Routing membership changes through the
    consensus log is simple but introduces latency (the change only takes effect after
    3-chain commit).  Systems that need rapid reconfiguration might prefer a separate fast
    path.

  \item \textbf{No threshold signatures.}  QCs in our system carry $2f{+}1$ individual
    Ed25519 signatures, which grows linearly in $f$.  Threshold (or aggregate) signatures
    could compress a QC into a single short proof but add cryptographic complexity.  We
    defer this optimisation to future work.
\end{itemize}

\section{Safety and Liveness}

\textit{Safety.} The consensus core follows HotStuff's 3-chain commit rule without
modification, so the original safety proof~\cite{hotstuff2019} applies.  The WAL and
snapshots add durability but do not change when or how blocks are committed.  Dummy
proposals are regular blocks that carry no-op commands; they pass through the same voting
and QC pipeline, so they cannot weaken safety.  For dynamic membership, our TLA+
specification (Section~\ref{sec:formal_verification}) provides machine-checked verification.

\textit{Liveness.}  HotStuff guarantees liveness after Global Stabilisation Time (GST)
assuming an honest leader.  Our pacemaker extension adds liveness under low load by ensuring
the 3-chain always has enough blocks to make progress.  For dynamic membership, liveness
depends on the new configuration having enough honest replicas, which is enforced by the
$n' \ge 4 \wedge n' \bmod 3 = 1$ invariant.

\section{Future Work}

Several directions remain open:

\begin{itemize}
  \item \textbf{BLS / threshold signatures} to reduce QC size from $O(f)$ signatures to a
    single constant-size proof.
  \item \textbf{Pipelined HotStuff} --- overlapping multiple proposal rounds to improve
    throughput beyond 1~commit/round.
  \item \textbf{Disk-paged block tree} to handle long-running deployments with millions of
    committed commands.
  \item \textbf{Reputation-based leader rotation} to penalise repeatedly faulty or slow
    leaders and improve practical latency.
  \item \textbf{WASM plugin support} for sandboxed, untrusted state machine implementations
    delivered as WASM modules.
\end{itemize}
